\documentclass[12pt]{article} % Тип документа (стаття) та розмір шрифту (12pt)
\usepackage[T2A]{fontenc}      % Кодування шрифтів для відображення кирилиці
\usepackage[utf8]{inputenc}    % Кодування вхідного файлу (стандарт для укр.мови)
\usepackage[ukrainian]{babel}  % Підключення словників та правил переносу укр.мови
\usepackage{amsmath, amssymb}  % Пакети для математичних формул та символів
\usepackage{array}             % Пакет для покращеної роботи з таблицями

\title{Практична робота: Математичні формули} 
\author{Яцюк Вікторія}                          
\date{\today}                                

\begin{document}
\maketitle % Команда, яка фактично малює заголовок, автора та дату на сторінці 
 \section* {Квадратне рівняння}  

Квадратне рівняння має вигляд:

\begin{equation}
ax^2 + bx + c = 0, \quad a \neg 0
\end{equation}

Його розв'язки залежать від значення дискримінанта $D$. Нижче наведено узагальнюючу таблицю:

\begin{table}[h]
\centering
\renewcommand{\arraystretch}{3}
\begin{tabular}{|c|c|c|}
\hline
\multicolumn{3}{|c|}{$ax^2 + bx + c = 0, a \neq 0, b \neq 0, c \neq 0$}\\ \hline
\multicolumn{3}{|c|}{$D = b^2 - 4ac$}\\ \hline
$\begin{aligned}
    x_1 &= \frac{-b + \sqrt{D}}{2a} \\
    x_2 &= \frac{-b - \sqrt{D}}{2a}
\end{aligned}$ &
$x = -\dfrac{b}{2a}$ &
коренів немає \\ \hline
\end{tabular}
\end{table}
\section*{Приклад розв'яхання}
Розв'яжемо рівняння $2x^2 - 4x - 6 = 0$:
\begin{align*}
a &= 2, \quad b = -4, \quad c = -6\\
D &= (-4)^2 - 4 \cdot 2 \cdot (-6) = 16 + 48 = 64\\
x &= \frac{-(-4)\pm \sqrt{64}}{2 \cdot 2} \\
x &=\frac{4 \pm 8}{4} \\
x_1 &= \frac{12}{4} = 3, \quad x_2 &= \frac{-4}{4} = -1
\end{align*}

\end{document} % Кінець документа
